% Ad Atticum 13.21 (ad-atticum-13-21) --- English translation
% Translated by Claude (Anthropic) from the Latin (Perseus).
% Style guide: STYLE.md.

\ciceroLetterOpener{CICERO ATTICO salutem}

\ciceroSection{1} I had given to Hirtius a letter of considerable size which I had written most recently at the Tusculan villa. To the one you have just sent me I shall reply another time.

\ciceroSection{2} Now I prefer other business. What can I say about Torquatus until I get something from Dolabella? As soon as I do, you shall both know straight away. I was expecting couriers from him today or at the latest tomorrow; as soon as they arrive, they shall be sent on to you. I am waiting on Quintus. For when I set out from the Tusculan villa on the eighth day before the Kalends, as you know, I dispatched couriers to him.

\ciceroSection{3} Now, to come back to the matter at hand: that \textit{inhibere} of yours which had so taken my fancy is something I now thoroughly dislike. For the word is altogether nautical. I knew that much, of course, but I had supposed that the oars were held steady when the rowers were ordered \textit{inhibere}. That it is not so I learned yesterday, when a ship was being put in at our villa. They do not hold the oars steady --- they row, but in a different fashion. That is the very opposite of \textit{epoch\=e} \textit{[Greek: epoch\=ei]}. So please see to it that in the book it stands as it did before. You will say the same to Varro, in case he has happened to make the change. And nothing is better than Lucilius's line: \textit{sustineas currum ut bonus saepe agitator equosque} --- ``hold back the chariot and the horses, as a good driver often does.'' And Carneades always compares the boxer's guard \textit{[Greek: probol\=en]} and the charioteer's reining-in to \textit{epoch\=e} \textit{[Greek: epoch\=ei]}. The rowers' \textit{inhibitio}, on the other hand, involves motion --- and indeed a more vigorous motion of rowing, that turns the ship toward its stern. You see how much more carefully I attend to this than to the gossip or to Pollio.

\ciceroSection{4} About Pansa too, if you have anything more definite (for I take it the thing has been made public); about Critonius, if there is anything \dag{}\textit{esset certe ne}\dag{} about Metellus and Balbinus.
